\section*{Chapter \ref{ch:m3:blockchains-for-traders} --- Blockchains for Traders}

\begin{solution}{exo:m3:blockchains-for-traders:1}
$20 \times (1 + \tfrac18 \times 0.25/0.5) = 21.25$ gwei.
\end{solution}

\begin{solution}{exo:m3:blockchains-for-traders:2}
About an hour (six blocks of ten minutes on average); about 12.8 minutes (two epochs) for Ethereum \omterm{def:m3:blockchains-for-traders:finality}{finality}.
\end{solution}

\begin{solution}{exo:m3:blockchains-for-traders:3}
The peg is held by arbitrage between redemption at one dollar and the market price; only holders with primary
access can do that arbitrage. If few can redeem, or redemption is closed (banking hours, a frozen reserve), the
market price can fall below a dollar with nothing to pull it back.
\end{solution}

\begin{solution}{exo:m3:blockchains-for-traders:4}
Full blocks: $\ln 2/\ln 1.125 = 5.9$, so six more than double it ($1.125^6 = 2.03$). Empty blocks:
$\ln 0.5/\ln 0.875 = 5.2$, so five leave 51\% and six leave 45\%.
\end{solution}

\begin{solution}{exo:m3:blockchains-for-traders:5}
$150\,000 \times 32 \times 10^{-9} \times 3\,000 = \$14.40$.
\end{solution}

\begin{solution}{exo:m3:blockchains-for-traders:6}
A \omterm{def:m3:blockchains-for-traders:rollup}{rollup}'s \omterm{def:m3:blockchains-for-traders:rollup}{sequencer} includes it at once and cheaply, but the base layer guarantees it only when the batch is
posted and (for an optimistic \omterm{def:m3:blockchains-for-traders:rollup}{rollup}) its challenge period has passed, or (for a validity \omterm{def:m3:blockchains-for-traders:rollup}{rollup}) its proof is
verified; until then the \omterm{def:m3:blockchains-for-traders:rollup}{sequencer}'s ordering could in principle be replaced.
\end{solution}

\begin{solution}{exo:m3:blockchains-for-traders:7}
The \omterm{def:m3:blockchains-for-traders:gas}{base fee} peaks at 39.20 gwei; blocks stop being full at block 11, once the \omterm{def:m3:blockchains-for-traders:gas}{base fee} has risen enough to bring
demand below the limit.
\end{solution}

\begin{solution}{exo:m3:blockchains-for-traders:8}
The peg also needs redemption to work when holders want it: banking hours, a custodian failure (as at Silicon
Valley Bank in 2023), a legal freeze or the issuer's own solvency can stop redemptions, and secondary prices then
fall regardless of what the reserves are worth.
\end{solution}

\begin{solution}{pb:m3:blockchains-for-traders:1}
\textbf{1.} \$5.40. \textbf{2.} 11.25 gwei after one, 105.45 gwei after twenty. \textbf{3.} \$48.35.
\textbf{4.} About four minutes (twenty 12-second slots). \textbf{5.} The \omterm{def:m3:blockchains-for-traders:gas}{base fee} is burned; the tip goes to the
block's proposer. \textbf{6.} $60/(150\,000 \times 10^{-9} \times 3\,000) - 2 = 131.33$ gwei. \textbf{7.} After
22 full blocks ($\ln 13.13/\ln 1.125 = 21.9$). \textbf{8.} The break-even rises to 264.67 gwei, reached after 28
full blocks. \textbf{9.} Competition for inclusion: other searchers bid for the same opportunity, and the tip, not
the \omterm{def:m3:blockchains-for-traders:gas}{base fee}, decides the order. \textbf{10.} It falls by up to 12.5\% a block until demand meets the target.
\textbf{11.} Yes, but its fills are final only when batches settle, liquidity and opportunities differ, and its
\omterm{def:m3:blockchains-for-traders:rollup}{sequencer} controls ordering. \textbf{12.} It hides the transaction from other searchers until inclusion, at the
cost of paying the builder rather than competing in the open \omterm{def:m3:blockchains-for-traders:mempool}{mempool}. \textbf{13.} So that proposers cannot profit
from inflating it, and to make fees predictable; burning also returns value to all holders. \textbf{14.} Their
transactions cost ten times more or wait. \textbf{15.} Which transactions and contracts are competing, and the tips
others pay. \textbf{16.} It prices congestion transparently and predictably, but it allocates space to the highest
payers, which in a spike are those with the most to extract. \textbf{17.} How long transfers take to be final and
what they cost in congestion, so that a cross-venue trade is priced with both. \textbf{18.} \omterm{def:m3:blockchains-for-traders:gas}{Gas} is a fixed cost per
transaction, so it takes a larger share of a small profit. \textbf{19.} \emph{Named result:} twenty full blocks
multiply the \omterm{def:m3:blockchains-for-traders:gas}{base fee} by 10.5 (10 to 105.45 gwei); the \$60 arbitrage stops paying after 22 full blocks.
\textbf{20.} Position and inclusion in a block.
\end{solution}

\begin{solution}{iq:m3:blockchains-for-traders:1}
The point after which a transaction cannot be reversed. A venue credits deposits only after it, so a transfer
between exchanges is locked for that time and exposed to prices; a reversed deposit is a loss for the venue.
\iqlookfor{settlement latency as inventory risk.}
\end{solution}

\begin{solution}{iq:m3:blockchains-for-traders:2}
The \omterm{def:m3:blockchains-for-traders:gas}{base fee} is set by protocol from the previous block's \omterm{def:m3:blockchains-for-traders:gas}{gas} used (up or down by at most 12.5\%) and burned; the
\omterm{def:m3:blockchains-for-traders:gas}{priority fee} is chosen by the user and paid to the proposer to be included before others.
\iqlookfor{the update rule, the burn, and the tip as the auction.}
\end{solution}

\begin{solution}{iq:m3:blockchains-for-traders:3}
Primary-market arbitrage by institutions that mint and redeem at one dollar against secondary prices. It fails if
reserves lose value or become inaccessible, if redemption is suspended or restricted, or if the issuer's solvency or
legal standing is doubted.
\iqlookfor{mechanism plus the conditions it needs.}
\end{solution}

\begin{solution}{iq:m3:blockchains-for-traders:4}
The bridged \omterm{def:m3:blockchains-for-traders:key}{token} is a claim on assets locked in the \omterm{def:m3:blockchains-for-traders:bridge}{bridge}: smart-contract bugs, compromised signing keys, the
\omterm{def:m3:blockchains-for-traders:bridge}{bridge}'s governance, and \omterm{def:m3:blockchains-for-traders:stable}{depegs} between the bridged and the native \omterm{def:m3:blockchains-for-traders:key}{token}.
\iqlookfor{the \omterm{def:m3:blockchains-for-traders:bridge}{bridge} as a concentrated custodian.}
\end{solution}

\begin{solution}{iq:m3:blockchains-for-traders:5}
Credit only after the chain's \omterm{def:m3:blockchains-for-traders:finality}{finality} (or a confirmation count set by the amount at risk and the chain's reorg
history); until then show it as pending. Risks: a reorganisation drops the transaction, a double spend, or a
fork.
\iqlookfor{risk-based confirmation policy.}
\end{solution}

\begin{solution}{iq:m3:blockchains-for-traders:6}
Observe the \omterm{def:m3:blockchains-for-traders:mempool}{mempool} and recent blocks: the distribution of tips of included transactions by block, the pending
queue and its tips, and the \omterm{def:m3:blockchains-for-traders:gas}{base fee} path. Estimate, for each tip, the probability of inclusion in the next two
blocks (\omterm{def:m3:blockchains-for-traders:gas}{base fee} projected, competition from pending transactions and expected new arrivals), and choose the lowest
tip above 95\%; update each block and back-test against realised inclusions.
\iqlookfor{data, a probabilistic model, and calibration against outcomes.}
\end{solution}
